Catorce trabajos, seis ejes. Las fichas están en el Anexo~B; aquí solo el recorte que usa el
plan. El comentario por paper está en §\ref{sec:estado-arte}.

\begin{enumerate}
  \item \textbf{Madurez de OpenAPI.} En 47 artículos (2011--2020), Casas et al.\ (2021) tratan
        OpenAPI/Swagger como el estándar de facto de las APIs REST: generación, pruebas y
        gobernanza. Coherencia con un modelo de negocio bancario no entra en esa revisión.
  \item \textbf{Correspondencia.} Shvaiko y Euzenat (2005) ordenan el \textit{schema matching}.
        El Bayadh et al.\ (2015), Chandrasekaran y Mago (2022) y Santos Sousa et al.\ (2026)
        cubren alineación ontológica y similitud. Llegan por separado. El score único del par
        contrato--SD se define en este plan, no en esos papers.
  \item \textbf{Dominio e interfaz.} Schmidt y Thiry (2020) revisan 27 estudios de microservicios
        con MDE y DDD. Sirve para no tratar el YAML como si fuera el \ServiceDomain{}.
  \item \textbf{Lo que no es el foco.} Testing, seguridad y selección de servicios (E5 y E6) usan
        OpenAPI para cobertura, vulnerabilidades o ranking por QoS. Eso se excluye en
        \Cref{sec:alcances}.
\end{enumerate}
